\chapter{Негауссова пространственная жёсткость и однородная долина потери ранга}

Детерминантная ветвь закрыта. Том V, однако, построил пространственную
Bott-суспензию и показал, что сумма квадратичного и скирмовского
производных членов может удерживать конечный радиус топологической
конфигурации. Проверим, способен ли тот же механизм остановить уход
семейного состояния к проекторной границе.

\section{Два разных вопроса устойчивости}

Необходимо различать:
\begin{enumerate}
 \item выбор однородной фазы, то есть существование конечного минимума по
 $R$ и мосту $B$;
 \item устойчивость пространственного размера уже выбранного
 топологического дефекта.
\end{enumerate}
Баланс Деррика относится ко второму вопросу. Если
\begin{equation}
 E(L)=aL+\frac bL,
 \qquad a,b>0,
\end{equation}
то радиус стабилизируется при
\begin{equation}
 L_*=\sqrt{b/a}.
 \label{eq:v6-spatial-stiffness-derrick-radius}
\end{equation}
Такая структура стандартна для моделей Скирма и Фаддеева--Скирма; см.
\path{arXiv:hep-th/9811176} и \path{arXiv:math-ph/0211010}.

Но производные члены не выбирают значение постоянного поля. В теориях
Ландау--де Жена упругая энергия штрафует пространственные искажения, а
однородную фазу выбирает отдельный объёмный потенциал. Это разделение
сохраняется и в проектной постановке.

\section{Однородная долина переживает все производные члены}

На чистой страте положим
\begin{equation}
 R=P=\operatorname{diag}(1,0,0),
 \qquad
 B_C=\operatorname{diag}(1,C),
 \qquad C\in M_2(\mathbb R).
 \label{eq:v6-spatial-stiffness-uniform-valley}
\end{equation}
Предыдущая непертурбативная проверка доказала
\begin{equation}
 \mathcal S_P(B_C)=0
\end{equation}
для любого $C$.

Теперь сделаем $R(x)=P$ и $B(x)=B_C$ постоянными в пространстве. Тогда
\begin{equation}
 dR=0,
 \qquad dB=0,
 \qquad P(dP)^2=0.
\end{equation}
Поэтому одновременно обращаются в нуль:
\begin{align}
 E_2[R]&=\int|dR|^2,\\
 E_2[B]&=\int|dB|^2,\\
 E_4[R]&=\int\Tr\bigl(P(dP)^2\bigr)^2,
\end{align}
а также любой локальный полином, содержащий хотя бы одну пространственную
производную.

\begin{theorem}[Производная жёсткость не насыщает однородный переход]
Никакая комбинация квадратичных, скирмовских или более высоких локальных
производных членов не поднимает постоянное семейство
\eqref{eq:v6-spatial-stiffness-uniform-valley}. Четырёхмерная долина
$C\in M_2(\mathbb R)$ остаётся точной нулевой долиной.
\end{theorem}

\section{Дискретный контроль}

На связном пространственном графе добавим положительную жёсткость
\begin{equation}
 E_{\rm grad}
 =\kappa\sum_{\langle xy\rangle}\|C_x-C_y\|^2.
\end{equation}
Её гессиан равен графовому лапласиану, тензорно умноженному на четыре
компоненты $M_2(\mathbb R)$. Для связного графа лапласиан имеет одну
постоянную нулевую моду, поэтому полный гессиан имеет ровно четыре нуля.
Все неоднородные уходы подавляются, но глобальный сдвиг
\begin{equation}
 C_x=C
 \qquad\text{для всех }x
\end{equation}
остаётся бесплатным.

Машинный аудит восьмиузлового кольца получил четыре нулевые моды и
положительный остальной спектр. На тридцати случайных матрицах $C$ полная
локальная и производная энергия постоянной конфигурации равна нулю с
машинной точностью.

\section{Что сохраняется от результатов Тома V}

Результаты Тома V не отменяются:
\begin{enumerate}
 \item относительная Bott-суспензия по-прежнему несёт заряды $+15/-15$;
 \item если проекторная фаза уже выбрана и заданы $a,b>0$, баланс
 Деррика может удерживать конечный радиус;
 \item проекторная кривизна даёт математически допустимый скирмовский
 инвариант.
\end{enumerate}

Но ранее уже доказано, что общий коэффициент $b/a$ из текущего следа не
выведен, а фермионный детерминант не гарантирует единственный положительный
четырёхпроизводный сектор. Новый результат сильнее в другом направлении:
даже идеально выбранные положительные $a$ и $b$ не решают однородную
потерю ранга, потому что обе производные энергии на ней равны нулю.

\begin{proposition}[Жёсткость удерживает частицу, но не рождает фазу]
Скирмовский механизм может стабилизировать размер уже существующего
пространственного дефекта, однако не способен сам выбрать проекторное
вакуумное многообразие и сделать интеграл по мосту коэрцитивным.
\end{proposition}

\section{Граница и следующий маршрут}

Для насыщения требуется локальный непроизводный член, который видит
поперечный блок $C$ даже при постоянных $R$ и $B$. Он не обязан быть
логарифмическим детерминантом, но должен оставаться положительным на
ядре $R$.

Том V содержит ещё один неиспользованный структурный намёк: модулярное
соответствие различает левое состояние и коммутантное действие. Поэтому
следующая проверка должна проверить, создаёт ли обязательная пара прямого и
противоположного весов локальную коэрцитивность на
$\ker R$ без ручного добавления дополнения $I-R$.

\begin{nogocriterion}[Запрет подмены рождения стабилизацией]
Нельзя считать конечный радиус решения $E_2+E_4$ доказательством рождения
материи. Сначала должен существовать конечный однородный минимум полного
локального функционала; только затем разрешён баланс пространственного
размера внутри ненулевого топологического сектора.
\end{nogocriterion}

\begin{protocol}
\begin{enumerate}
 \item проекторный скирмовский инвариант:
 \textbf{математически допустим};
 \item общий коэффициент $E_2/E_4$:
 \textbf{не выведен};
 \item баланс Деррика при $a,b>0$:
 \textbf{стабилизирует радиус};
 \item локальная энергия постоянной долины $B_C$:
 \textbf{ноль};
 \item квадратичная пространственная жёсткость на долине:
 \textbf{ноль};
 \item скирмовская энергия на долине:
 \textbf{ноль};
 \item число глобальных нулевых направлений после жёсткости:
 \textbf{четыре};
 \item насыщение однородного перехода:
 \textbf{нет};
 \item стабилизация размера уже созданного дефекта:
 \textbf{условно возможна};
 \item следующая проверка:
 \textbf{модулярная двойственная коэрцитивность моста};
 \item рождение материи:
 \textbf{не закрыто}.
\end{enumerate}
\end{protocol}

Машинный сертификат сохранён в
\path{s2t_v6_nongaussian_spatial_stiffness_saturation_gate_results.json}.